\section{The Batch That Fixes It}
\label{sec:ch03-the-batch-that-fixes-it}

\index{DataLoader!source-generated|(}%
A DataLoader in Hot Chocolate is a method with an attribute on it. Here is the
whole of \code{SpeakerDataLoader.cs}:

\begin{minted}{csharp}
using Microsoft.EntityFrameworkCore;

namespace Sessions;

public static class SpeakerDataLoader
{
    /// <summary>
    /// Every speaker a batch of sessions asked for, keyed by identifier.
    /// </summary>
    [DataLoader]
    public static async Task<Dictionary<int, Speaker>>
        GetSpeakerByIdAsync(
            IReadOnlyList<int> ids,
            ConferenceContext db,
            CancellationToken ct)
        => await db.Speakers
            .Where(s => ids.Contains(s.Id))
            .ToDictionaryAsync(s => s.Id, ct);
}
\end{minted}

\index{DataLoader!signature}%
Every part of that signature is doing something. It takes a list of keys and
returns a dictionary from key to value, which is how the engine hands over the
whole batch and gets back something it can take apart again.
\code{ConferenceContext} is an ordinary service parameter, resolved the same
way it is on a resolver. And the method is a plain static function: it does not
know what a session is, it is not attached to a type, and it would be just as
happy loading speakers for something else.

\index{DataLoader!generated name}%
The name is not decoration either. Strip \code{Get} from the front and
\code{Async} from the end of \code{GetSpeakerByIdAsync}, and what is left is
what the generator calls the thing it writes. Below are the two declarations it
emits, lifted out of their namespace and separated from the doc comment and the
fetch body that sit between them in the real file:

\begin{minted}[fontsize=\footnotesize]{csharp}
public interface ISpeakerByIdDataLoader
    : global::GreenDonut.IDataLoader<int, global::Sessions.Speaker>
{
}

public sealed partial class SpeakerByIdDataLoader
    : global::GreenDonut.DataLoaderBase<int, global::Sessions.Speaker>
    , ISpeakerByIdDataLoader
\end{minted}

That is the same naming rule chapter~\ref{ch:hot-chocolate-zero} met on
\code{AddSessionsTypes}, applied to a different input, and it has the same
consequence: the type you are about to reference does not exist until you
build, and its name came from a method name rather than from anything you
declared.

\index{resolver!consuming a DataLoader}%
\code{SessionType.cs} changes by one parameter and one line of body. The
context goes, the loader arrives, and the query becomes a load:

\begin{minted}[highlightlines={11,13}]{csharp}
namespace Sessions;

[ObjectType<Session>]
public static partial class SessionType
{
    /// <summary>
    /// The person giving this session.
    /// </summary>
    public static async Task<Speaker?> GetSpeakerAsync(
        [Parent] Session session,
        ISpeakerByIdDataLoader speakerById,
        CancellationToken ct)
        => await speakerById.LoadAsync(session.SpeakerId, ct);
}
\end{minted}

\index{using directive!not needed for a DataLoader}%
The \code{using Microsoft.EntityFrameworkCore;} at the top is gone, because
this file no longer knows there is a database. That is worth more than it
looks: the resolver's job is now to say which speaker it wants, and where
speakers come from is somewhere else entirely. When
chapter~\ref{ch:first-subgraph} turns this service into a subgraph, that
separation is what lets the same resolver shape answer a router.

\index{AddSessionsTypes@\code{AddSessionsTypes}!registers the DataLoader}%
Now the part the documentation and the compiler disagree about.
\code{Program.cs} does not change. The documentation says it should: it
describes registration as declaring a module for your assembly, after which
\enquote{the source generator then emits a registration extension method named
after the module}~\autocite{chillicreamdataloader}, which you then call. Do
none of that, build, and read what the generator wrote into the method
chapter~\ref{ch:hot-chocolate-zero} already calls. Its enclosing namespace and
class are dropped here, and three statements that did not change are replaced
by the comments naming them:

\begin{minted}[fontsize=\footnotesize, breakanywhere]{csharp}
public static IRequestExecutorBuilder AddSessionsTypes(this IRequestExecutorBuilder builder)
{
    builder.AddTypeExtension(typeof(global::Sessions.Query));
    // the SessionType descriptor registration
    builder.AddDataLoader<global::Sessions.ISpeakerByIdDataLoader, global::Sessions.SpeakerByIdDataLoader>();
    // the root Query type, and the Session object type
    return builder;
}
\end{minted}

\code{AddSessionsTypes} registered it, unasked. The attribute is the whole of
the wiring, exactly as \code{[QueryType]} was in
chapter~\ref{ch:hot-chocolate-zero}, and the module attribute is something you
would need if you wanted the registrations somewhere else. I would rather find
this out from the generated file than from a page, which is the argument for
reading generated code at least once per library.

\index{statement count!with a DataLoader}%
Delete the database, restart, send the same request as
section~\ref{sec:ch03-what-one-query-costs}:

\begin{minted}[fontsize=\footnotesize, breakanywhere]{text}
-- statement 1
SELECT "s"."Id", "s"."Abstract", "s"."DurationMinutes", "s"."SpeakerId", "s"."StartsAt", "s"."Title"
FROM "Sessions" AS "s"
ORDER BY "s"."StartsAt"
-- statement 2
SELECT "s"."Id", "s"."Bio", "s"."Name"
FROM "Speakers" AS "s"
WHERE "s"."Id" IN (@ids1, @ids2, @ids3)
\end{minted}

\index{DataLoader!deduplicates keys}%
Two statements instead of five, and the interesting number is not the two. It
is the three parameters in that \code{IN} list, for four sessions. Ada Fischer
gives two of them, both calls asked for speaker~1, and the loader asked SQLite
once. Nothing in \code{SessionType.cs} knows that happened.

\index{DataLoader!missing key resolves to null}%
One behavior to know before you rely on it. The dictionary you return does not
have to contain every key you were given, and a key you leave out resolves to
null rather than to an error. That is the right default for
\code{speaker: Speaker}, a field the schema already permits to be null. It is
the wrong default to discover by accident on a non-null field, and
chapter~\ref{ch:null-blast-radius} is about what happens then.

The count has stopped tracking the size of the list. A longer schedule makes
the second statement's \code{IN} list longer rather than making a third
statement, which is the property the chapter was after and the one
chapter~\ref{ch:router-n-plus-1} has to rebuild from scratch once the speakers
live in another service.
\index{DataLoader!source-generated|)}%
